% ID: SEL-D0055-OA-E02
% ESTADO: validada. S056.
\textbf{a)} Las ramas son suaves en sus interiores. En $2$, ambos límites coinciden con $f(2)=2$. En $4$, ambos coinciden con $f(4)=1$. Por tanto, $f$ es continua en todo $\mathbb R$. Las derivadas de las ramas son $-1$, $-4/x^2$ y $2x-4$. En $2$ las dos derivadas laterales valen $-1$, de modo que $f'(2)=-1$. En $4$ valen $-1/4$ y $4$, así que $f$ sólo deja de ser derivable en $4$.

\textbf{b)} Los dos primeros tramos decrecen y el último crece. La continuidad enlaza los dos descensos en $2$. El único extremo local es un mínimo en $(4,1)$, que también es absoluto.

\textbf{c)} $f(3)=4/3$ y $f'(3)=-4/9$. La tangente es
\[
y-\tfrac43=-\tfrac49(x-3)
\quad\Longleftrightarrow\quad
\boxed{y=-\tfrac49x+\tfrac83}.
\]
